\chapter{Mơ và Tỉnh}
\markboth{Mơ và Tỉnh}{Dreaming and Awakening}

\section{Mơ lớn để đời không bị nhỏ lại.}
\begin{truyen}{Mơ lớn để đời không bị nhỏ lại.}{Dreams can protect dignity from routine.}
\chuhoa{C}{ậu học sinh ở vùng quê nói mình muốn làm kỹ sư dù quanh em ít ai tin. Ước mơ ấy chưa cho em kết quả ngay, nhưng nó giữ em khỏi buông xuôi theo một đời sống quá dễ dãi. Có những giấc mơ không phải để khoe, mà để giữ cho con người không sống thấp hơn khả năng của mình.}
\textit{(A rural student said he wanted to become an engineer, though few around him believed it. The dream did not produce immediate results, but it kept him from surrendering to an overly passive life. Some dreams are not for display; they keep a person from living beneath their potential.)}
\end{truyen}
\begin{baihoc}
Mơ đúng cách làm con người có phương hướng và lòng tự trọng.
\textit{(Dreaming well gives direction and self-respect.)}
\end{baihoc}
\begin{ghinhoanh}
\textbf{Short English to remember:}

Dreams can protect your direction.
\end{ghinhoanh}
\begin{tuhocnhanh}
1. Ước mơ nào đang giữ em khỏi sống hời hợt? \textit{(Which dream is keeping you from living superficially?)}

2. Em đã làm gì để giấc mơ ấy đi gần hơn với thực tế? \textit{(What have you done to bring that dream closer to reality?)}
\end{tuhocnhanh}
\ngancach

\section{Mơ đẹp nhưng quên tỉnh để làm việc thực tế.}
\begin{truyen}{Mơ đẹp nhưng quên tỉnh để làm việc thực tế.}{A dream without effort stays decorative.}
\chuhoa{A}{nh nói rất hay về tương lai, về thương hiệu riêng, về những bước ngoặt lớn. Nhưng ngày nào cũng trôi qua mà không có giờ học, không có luyện tập, không có kỷ luật. Giấc mơ nếu không đi cùng sự tỉnh táo của lao động thì chỉ là một tấm poster treo trong đầu.}
\textit{(He spoke beautifully about the future, personal branding, and major breakthroughs. Yet each day passed without study time, practice, or discipline. When a dream is not joined with the awake reality of work, it becomes only a poster hanging in the mind.)}
\end{truyen}
\begin{baihoc}
Mơ không sai; sai là dùng giấc mơ để thay cho hành động.
\textit{(Dreaming is not wrong; what is wrong is using the dream in place of action.)}
\end{baihoc}
\begin{ghinhoanh}
\textbf{Short English to remember:}

A dream needs discipline.
\end{ghinhoanh}
\begin{tuhocnhanh}
1. Giấc mơ nào của em đang thiếu hành động cụ thể? \textit{(Which of your dreams is missing concrete action?)}

2. Em có đang thích tưởng tượng hơn là bắt tay vào làm không? \textit{(Do you enjoy imagining more than actually working?)}
\end{tuhocnhanh}
\ngancach

\section{Tỉnh ra sau một ảo tưởng về bản thân.}
\begin{truyen}{Tỉnh ra sau một ảo tưởng về bản thân.}{Awakening can be painful but necessary.}
\chuhoa{E}{m từng nghĩ mình học ổn vì chưa gặp bài khó thật. Khi điểm số sa xuống, em sốc và thấy bản thân như bị kéo khỏi một giấc mơ dễ chịu. Nhưng chính cú tỉnh ấy làm em bắt đầu học lại nghiêm túc hơn.}
\textit{(The student once thought the academic level was fine because true difficulty had not yet arrived. When the scores dropped, the shock felt like being pulled out of a comfortable dream. Yet that awakening became the beginning of more serious learning.)}
\end{truyen}
\begin{baihoc}
Có những cú tỉnh làm đau lòng tự ái, nhưng cứu tương lai.
\textit{(Some awakenings hurt pride, but save the future.)}
\end{baihoc}
\begin{ghinhoanh}
\textbf{Short English to remember:}

A painful awakening can save you.
\end{ghinhoanh}
\begin{tuhocnhanh}
1. Em đang ảo tưởng ở điểm nào về bản thân? \textit{(In what area might you be holding an illusion about yourself?)}

2. Dấu hiệu thực tế nào em nên nhìn thẳng hơn? \textit{(Which real signs should you face more honestly?)}
\end{tuhocnhanh}
\ngancach

\section{Vẫn giữ khả năng mơ dù đời không còn dễ.}
\begin{truyen}{Vẫn giữ khả năng mơ dù đời không còn dễ.}{Maturity should not kill wonder.}
\chuhoa{N}{gười lớn thường dễ cười vào những điều trong trẻo của người trẻ. Nhưng có một cô giáo dù đã đi qua nhiều va vấp vẫn giữ thói quen khuyến khích học sinh tưởng tượng, đặt câu hỏi và tin rằng điều tốt có thể được xây nên. Tỉnh táo không nhất thiết phải khô cằn.}
\textit{(Adults often laugh at the innocence of young dreams. Yet one teacher, after many hardships, still encouraged students to imagine, to question, and to believe that good things can be built. Being awake does not require becoming dry.)}
\end{truyen}
\begin{baihoc}
Trưởng thành đẹp là vừa tỉnh vừa còn biết mơ.
\textit{(Beautiful maturity is being awake while still able to dream.)}
\end{baihoc}
\begin{ghinhoanh}
\textbf{Short English to remember:}

Stay awake, but keep wonder.
\end{ghinhoanh}
\begin{tuhocnhanh}
1. Em có đang đánh mất khả năng mơ vì sợ thất vọng không? \textit{(Are you losing the ability to dream because you fear disappointment?)}

2. Em muốn giữ lại nét trong trẻo nào khi mình lớn lên? \textit{(What kind of innocence do you want to preserve as you grow up?)}
\end{tuhocnhanh}
\ngancach

\section{Tỉnh để chỉnh giấc mơ, không phải giết giấc mơ.}
\begin{truyen}{Tỉnh để chỉnh giấc mơ, không phải giết giấc mơ.}{Awareness should refine dreams, not bury them.}
\chuhoa{K}{hi nhận ra mình không hợp với ngành từng theo đuổi, cậu sinh viên không coi đó là hết. Cậu chỉ chuyển hướng, giữ tinh thần phục vụ con người nhưng đổi con đường cho phù hợp hơn. Tỉnh không phải là phản bội giấc mơ; nhiều khi là cứu nó khỏi hình thức sai.}
\textit{(When he realized he did not fit the field he had once pursued, the student did not see it as the end. He simply changed direction, keeping the desire to serve people while choosing a more fitting path. Awakening is not betraying a dream; sometimes it rescues the dream from the wrong form.)}
\end{truyen}
\begin{baihoc}
Giấc mơ trưởng thành là giấc mơ biết điều chỉnh để còn sống được lâu.
\textit{(A mature dream is one that knows how to adjust in order to live longer.)}
\end{baihoc}
\begin{ghinhoanh}
\textbf{Short English to remember:}

Refine the dream, do not kill it.
\end{ghinhoanh}
\begin{tuhocnhanh}
1. Giấc mơ nào của em cần được chỉnh lại cho thực hơn? \textit{(Which dream of yours needs refining to become more realistic?)}

2. Điều cốt lõi nào trong giấc mơ ấy em muốn giữ lại? \textit{(What core element of that dream do you want to preserve?)}
\end{tuhocnhanh}
\ngancach
